\documentclass{situationist-preprint}

\title{Descartes, Husserl and the Cogito}
\msid{PHE-20-b}
\msfamily{PHE-20}
\msdiscipline{Phenomenology}
\mstype{Companion manuscript}
\mskeywords{cogito; Descartes; Husserl; transcendental ego; thinking substance; residue}

\begin{document}
\maketitle

\begin{abstract}
Descartes's cogito and Husserl's reduction both end in an ego that survives universal doubt or bracketing. Husserl held that Descartes, having reached the ego, mistook it for a thinking substance, a piece of the world. We show that PHE-20-a vindicates this verdict precisely. The certainty of the cogito is the certainty of the residue of Theorem~3.1 of PHE-20-a, and Descartes's own restriction, that I exist as long as I think, states the residue's mode of presence. The inference to a substance that persists across acts is blocked by Theorem~3.2 of PHE-20-a.
\end{abstract}

\section{Two survivors of doubt}

In the second \emph{Meditation}, Descartes finds that whatever he doubts, he cannot doubt that he is, while he thinks; and he adds that he is certain of his existence only for as long as he is thinking \citep{Descartes1641}. He then asks what this \enquote{I} is and answers that it is a thing that thinks, which the sixth \emph{Meditation} treats as a substance distinct from the body and persisting through its successive thoughts. Husserl took the first step as the model of the reduction and rejected the second: Descartes, he argued, made the ego a \emph{substantia cogitans}, a remnant of the world left over after doubt, and so failed to make the transcendental turn \citep[§10]{Husserl1950}.

\section{The cogito as residue}

\begin{proposition}\label{p:certain}
The certainty of the cogito is the certainty of the residue: in each act of doubt, the act itself is not doubted by it.
\end{proposition}

\begin{proof}
Universal doubt is a case of bracketing in the sense of Definition~2.1 of PHE-20-a: the positing of each item as existing is put out of action. By Theorem~3.1 of PHE-20-a, every such act leaves its own act unbracketed. So the doubt cannot reach the act of doubting, which is lived through as performed, and in it the one who performs it.
\end{proof}

Descartes's restriction, that the certainty lasts as long as the thinking, is the temporal form of Proposition~\ref{p:certain}. The residue is present in each act, not before or after it (Corollary~3.3 of PHE-20-a), and the cogito accordingly certifies an existence coextensive with the act.

\begin{proposition}\label{p:substance}
The cogito does not establish a thinking substance that persists across acts.
\end{proposition}

\begin{proof}
A substance that persists across acts persists across acts in the sense of Assumption~2.4 of PHE-20-a, and so can be bracketed by some act. By Theorem~3.2 of PHE-20-a, nothing persisting survives every reduction. Since the certainty of the cogito is the certainty of what survives the doubt in which it is attained (Proposition~\ref{p:certain}), it does not extend to a persisting substance.
\end{proof}

Proposition~\ref{p:substance} is the precise form of Husserl's verdict. It also restates, in phenomenological terms, the point Kant made against the rational psychologist: from the \enquote{I think} that accompanies representations, nothing follows about a simple, persisting subject \citep[A341--405/B399--432]{Kant1781}. It does not follow that Sartre was right that the ego is merely a transcendent object \citep{Sartre1937}: what is excluded is the persisting ego, not the residue, which is present in every act.

\begin{objection}
Descartes needs the persisting substance for his later arguments, and the certainty of a momentary ego is too thin to build on.
\end{objection}

\begin{reply}
That is Husserl's point. The certainty of the cogito is thin in exactly the way the residue is: it is renewed in each act and guaranteed in none beyond it. What Descartes built on it required a further premise, and Proposition~\ref{p:substance} shows that the cogito does not supply it.
\end{reply}

\section{Conclusion}

The cogito and the reduction reach the same residue. Descartes was right about its certainty and its bond to the act, and wrong to make it a thing that outlasts the act. PHE-20-a explains both verdicts at once.

\begin{thebibliography}{9}\small
\bibitem[Descartes(1641)]{Descartes1641} Descartes, R. (1641). \emph{Meditationes de prima philosophia}. Michel Soly.
\bibitem[Husserl(1950)]{Husserl1950} Husserl, E. (1950). \emph{Cartesianische Meditationen und Pariser Vortr\"age}. Husserliana I. Ed. S. Strasser. Martinus Nijhoff.
\bibitem[Kant(1781)]{Kant1781} Kant, I. (1781/1787). \emph{Kritik der reinen Vernunft}. Hartknoch.
\bibitem[Sartre(1937)]{Sartre1937} Sartre, J.-P. (1937). La transcendance de l'ego. \emph{Recherches philosophiques}, 6.
\end{thebibliography}

\end{document}
